\begin{example}
    Рассмотрим $S \subset [0, 1]^{[0, 1]}$, состоящее из функций $x\colon [0, 1] \to [0, 1]$, таких что они совпадают с $0$ везде, кроме не более чем счётного множества точек.
    Оно не является топологическим компактом (оно незамкнуто -- а в хаусдорфовом пространстве компакт замкнут), но является секвенциальным компактом.

    Пусть $\{x_n\} \subset S$ -- произвольная последовательность.
    Пусть $M(x) = \{t \in [0, 1] \mid x(t) \neq 0\}$.
    При всяком $n \in \N$ множество $M(x_n)$ не более чем счётно, а значит и
    \[
        M = \bigcup_{n = 1}^{\infty} M(x_n),
    \]
    тоже не более чем счётно.
    Занумеруем его некоторым образом: $M = \{t_k\}_{k = 1}^{\infty}$.
    Мы можем выделить сходящуюся подпоследовательность из $x_n(t_1)$, некоторую $x_{n_m(1)}(t_1) \rightarrow y_1$.
    Потом из подпоследовательности выделить сходящуюся подпоследовательность $x_{n_m(2)}(t_2) \rightarrow y_2$, и так далее...

    Финальную подпоследовательность мы можем взять диагонально -- $x_{n_m(m)}$.
    Это действительно сработает, поскольку $\forall k \in \N$ последовательность $x_{n_m(m)}$ при $m \geq k$ является подпоследовательностью $x_{n_m(k)}$, а для неё верна
    \[
        x_{n_m(k)}(t_k) \rightarrow y_k.
    \]
    А для $t \in [0, 1] \setminus t_{k}$ все $\{x_n\}$ тождественно равны $0$.
\end{example}

\section{Аксиомы отделимости.}
\begin{definition}
    Пусть $(X, \tau) \in \cat{Top}$.
    Тогда
    \begin{enumerate}
        \item Говорят, что $X$ удовлетворяет первой аксиоме отделимости ($T1$-пространство) если $\forall x \in X \hookrightarrow \{x\}$ -- замкнуто.

        Это эквивалентно тому, что для всяких $x, y \in X\colon x \neq y \ \exists U(x) \not\ni y \wedge \exists V(y) \not\ni x$.
        \item Говорят, что $X$ удовлетворяет второй аксиоме отделимости (хаусдорфово пространство), если любые две точки можно отделить окрестностями друг от друга, то есть
        \[
            \forall x, y \in X\colon x \neq y \ \exists U(x) \wedge V(y) \hookrightarrow U \cap V = \emptyset.
        \]
    \end{enumerate}
\end{definition}
\begin{note}
    Очевидно, что в хаусдорфовом пространстве всякая последовательность имеет не более одного предела.
\end{note}
\begin{proposition}
    В хаусдорфовом пространстве $(X, \tau)$ компакт $K$ -- замкнут.
\end{proposition}
\begin{proof}
    Пусть $x \in X \setminus K$.
    Тогда для всякой точки $y \in K$ существуют $V_y(y)$ и $U_y(x)$ такие, что $V_y \cap U_y = \emptyset$.

    Заметим, что семейство $\{V_y\}_{y \in K}$ -- открытое покрытие $K$.
    Следовательно, из него можно выделить конечное подпокрытие $\{V_{y_k}\}_{k = 1}^{N}$.
    Но, тогда, $\bigcup_{k = 1}^{N} V_{y_k}$ и $\bigcap_{k = 1}^{N} U_{y_k}$ -- непересекающиеся окрестности $K$ и $x$ соответственно, и, следовательно, $\bigcap_{k = 1}^{N} U_{y_k}(x) \cap K = \emptyset$.
    Поэтому $x$ -- не точка прикосновения $K$ и, таким образом, $K$ -- замкнуто.
\end{proof}
\begin{proposition}
    В хаусдорфовом пространстве $(X, \tau)$ секвенциальный компакт $K$ является секвенциально замкнутым.
\end{proposition}
\begin{proof}
    Пусть $[K]_{seq} \subset X$ -- секвенциальное замыкание и $x \in [K]_{seq}$.
    Тогда $\exists \{x_n\} \subset K$ такая, что $x_n \xrightarrow{\tau} x$.
    Но она имеет сходящуюся подполседовательность, сходящуюся к элементу $K$.
    Поскольку последовательности в хаусдорфовом пространстве имеют не более одного предела, то $x \in K$ и $K$ -- секвенциально замкнуто.
\end{proof}
\begin{proposition}
    Пусть $(X, \tau)$ -- хаусдорфово пространство и $T$ -- непустое множество.
    Тогда $X^T$ -- хаусдорфово пространство.
\end{proposition}
\begin{proof}
    Пусть $x, y \in X^T\colon x \neq y$.
    Тогда $\exists t \in T\colon x(t) \neq y(t)$.
    В силу хаусдорфовости $X$ существуют $U(x(t))$ и $V(y(t))$ такие, что $U \cap V = \emptyset$.
    А значит и $\pi^{-1}_t(U) \cap \pi^{-1}_t(V) = \emptyset$, поэтому $\pi^{-1}_t(U)$ и $\pi^{-1}_t(V)$ являются искомыми непересекающимися окрестностями $x$ и $y$.
\end{proof}
\section{Счётная компактность.}
\begin{definition}
    Пусть $(X, \tau) \in \cat{Top}$.
    Множество $K \subset X$ называется счётным компактом, если из всякого счётного открытого покрытия $K$ можно выделить конечное подпокрытие.
\end{definition}
\begin{note}
    Очевидно, что всякий компакт является счётным компактом.
\end{note}
\begin{proposition}
    Если $(X, \tau) \in \cat{Top}$ и $K$ -- секвенциальный компакт, то $K$ -- счётный компакт.
\end{proposition}
\begin{proof}
    Пусть $K$ не является счётным компактом, то есть существует такое счётное покрытие $P = \{V_n\}_{n = 1}^{\infty}$, что из него нельзя выделить конечное подпокрытие.

    Тогда множества $W_n = \bigcup_{k = 1}^{n} V_k$ -- не являются покрытиями $K$ и $\forall n \in \N \  \exists x_n \in K \setminus W_n$.
    Последовательность $x_n$ имеет сходящуюся подпоследовательность $\{x_{n_k}\}\colon x_{n_k} \xrightarrow{\tau} x \in K$.
    Поскольку $P$ -- покрытие, то существует некоторое $j$ такое, что $V_j \ni x$.
    Но, тогда, поскольку $V_j$ -- открыто, то начиная с некоторого $N$ все элементы подпоследовательности должны попадать в неё -- чего не происходит по построению последовательности.

    Следовательно $K$ -- счётный компакт.
\end{proof}
\begin{proposition}
    Пусть $(X, \tau) \in \cat{Top}$ и удовлетворяет аксиоме счётности.
    Тогда, если $K \subset X$ -- счётный компакт, то $K$ -- секвенциальный компакт.
\end{proposition}
\begin{proof}
    Пусть $\{x_n\} \subset K$.

    Покажем, что существует точка $x \in K$ такая, что в любой её окрестности лежит бесконечно много элементов $\{x_n\}$ (где под <<бесконечно много>> понимается именно множество индексов элементов, попавших в окрестность).
    Пусть такой точки не существует.
    То есть $\forall x \in K$ существует $U(x)$ такая, что в ней лежит лишь конечное число элементов последовательности -- мощность $\{n \in \N \mid x_n \in U(x)\}$ конечна.

    Рассмотрим $M \subset \N$, такое что $|M| < +\infty$.
    Пусть
    \[
        S_M = \left\{x \in K \mid \{n \in \N \mid x_n \in U(x)\} = M \right\}.
    \]
    Подмножеств натуральных чисел, имеющих конечную мощность -- счётное число.
    Пусть теперь
    \[
        V_M = \bigcup_{x \in S_M} U(x) \in \tau.
    \]
    Такие множества являются не более чем счётным покрытием $K$.
    В силу счётной компактности $K$ оно имеет конечное подпокрытие $\{V_{M_k}\}_{k = 1}^{N}$.

    Пусть теперь $M = \bigcup_{k = 1}^{N} M_k$.
    Оно -- конечно.
    Но, по определению $S_M$ и $V_M$, если $n \notin M$, то оно не лежит и в $V_{M_k}$ при всяком $k$.
    Противоречие -- индексное множество последовательности бесконечно, а $M$ имеет конечную мощность.

    Значит существует точка $x \in K$ такая, что в любой её окрестности лежит бесконечно много элементов $\{x_n\}$.
    Поскольку $(X, \tau)$ удовлетворяет аксиоме счётности, то существует счётная убывающая локальная база $\{W_k\}_{k = 1}^{\infty}$.
    Пусть теперь $\{x_{n_k}\}$ выбирается так, чтобы $n_{k + 1} > n_{k}$ и $x_{n_k} \in W_k$.
    Такая подпоследовательность будет сходиться к $x \in K$ -- а следовательно $K$ является секвенциальным компактом.
\end{proof}
\section{Limit-point compactness.}
\begin{definition}
    Если $(X, \tau) \in \cat{Top}$ и $E \subset X$, то $x \in X$ называется предельной точкой $E$, если для всякой окрестности $U(x)$ верно, что $U(x) \setminus \{x\} \cap E \neq \emptyset$.
\end{definition}
\begin{definition}
    Будем говорить, что $K \subset (X, \tau)$ является limit-point compact\footnote{Я не знаю как это нормально на русском сказать...}, если всякое бесконечное $E \subset K$ имеет в $K$ предельную точку.
\end{definition}
\begin{proposition}
    Если $(X, \tau) \in \cat{Top}$ и $K \subset X$ -- счётный компакт, то $K$ и limit point compact.
\end{proposition}
\begin{proof}
    Пусть $K$ не обладает свойством limit point compactness.
    Тогда для всякого бесконечного $E \subset K$ никакая точка $K$ не является предельной, то есть
    \[
        \forall x \in K \ \exists U(x)\colon U(x) \setminus \{x\} \cap E = \emptyset.
    \]
    Без ограничения общности будем считать $E$ -- счётным, поскольку если для несчётного (или более мощного) $E$ это свойство не выполнено, то и для всякого его подмножества оно тем более невыполнено.
    Рассмотрим семейство $P = \{U(x) \mid x \in E\} \cup \left\{ \bigcup_{x \in K \setminus E} U(x) \right\}$.
    Заметим, что множество $W = \bigcup_{x \in K \setminus E} U(x)$ -- открыто и дизъюнктивно с $E$.
    Таким образом мы получили $P$ -- не более чем счётное покрытие $K$.
    Но оно не имеет конечного подпокрытия, поскольку всякая точка $x \in E$ лежит только в одном множестве $U(x)$.

    Противоречие.
    Следовательно, $K$ является limit point compact.

\end{proof}
